\scanpage{158}
\mspage{V.30}
% archive sheet number: 157

L'homomorphisme $u$, (\deleted{et} \inserted{ou ce qui revient au même,} le triangle (6.2.2) correspondant) domine l'étude cohomologique des quatre schémas $V_s\simeq V_{\bar s},V,V_{\unsure{\eta|y}},V_{\bar\eta}$, en ce sens qu'il permet de retrouver formellement la cohomologie (à coefficients constants) des schémas envisagés. Appliquant le foncteur ``oubli de $\unsure{\pi|\Pi}$'' à 6.2.2, on trouve le triangle exact
\diagram{\[
\tag*{(6.2.3)}
\begin{tikzcd}[ampersand replacement=\&,column sep=small]
\& {\mathbb{R}\unsure{\Gamma|\mathscr{T}}(V_s,\Lambda_{\unsure{\nu|v}})=L^\bullet}\arrow[dl] \& \\
{R^{\bullet\,\mathrm{oub}}=\mathbb{R}\Gamma(\bar V,\Lambda_{\nu\unsure{,|\,}V_{\bar\eta},\bar V})}\arrow[rr] \& \& {\mathbb{R}\Gamma(V_{\bar\eta},\Lambda_\nu)=K^{\bullet\,\mathrm{oub}}}\arrow[ul,"{\unsure{\kappa|k|h}}^*"]
\end{tikzcd}
\]}
qui donne la cohomologie (sans opérateurs) de $V_{\bar\eta}$ et $V_s$ et l'hom du premier dans le second défini par l'inclusion
\diagram{\[
\begin{tikzcd}[ampersand replacement=\&]
V_{\bar s}\arrow[r,"\bar i"]\arrow[rr,bend right=25,dashed,"\kappa"'] \& \bar V \& V_{\bar\eta}\arrow[l,"\sim"',"\mathrm{coh.}"]
\end{tikzcd}
\]}.
D'autre part, \unsure{appliquons|appliquant} le foncteur $\mathbb{R}\Gamma^\pi$ \inserted{à (6.2.1), (6.2.2)}, on trouve
\diagram{\[
\tag*{(6.2.4)}
\begin{tikzcd}[ampersand replacement=\&,column sep=small]
\& {\mathbb{R}\Gamma(V_s,\Lambda_\nu)\mathbin{\mathop{\otimes}^{\mathrm L}}_{\Lambda_\nu}\mathbb{R}\Gamma^\pi(\Lambda_\nu)\simeq\deleted{\unsure{\mathbb{R}\Gamma}}L^\bullet\oplus L^\bullet[-1]}\arrow[dl] \& \\
{\mathbb{R}\Gamma^\pi(R^\bullet)}\arrow[rr] \& \& {\mathbb{R}\Gamma(V_\eta,\Lambda_\nu)\simeq\mathbb{R}\Gamma^\pi(K^\bullet)}\arrow[ul,"\mathbb{R}\Gamma^\pi(\kappa^*)\unsure{=|\sim}"]
\end{tikzcd}
\]}
\inserted{\unsure{dégager la formule|dégén. formule}}
% margin note written slanted
\margin{\unsure{remonter|remettre} \unsure{ce|le} diagramme et le commenter !
\diagram{\begin{tikzcd}[ampersand replacement=\&]
V \& V_\eta\arrow[l,"\deleted{\unsure{g}}\unsure{j_s|j_\eta|j}"]\\
{V_s\simeq V_{\bar s}}\arrow[u,"\deleted{\unsure{h_s}}i"]\arrow[r,dashed,"\kappa"] \& {\bar V\overset{\deleted{\unsure{\mathrm{coh.}}}}\simeq V_{\bar\eta}}\arrow[u,"h_\eta\deleted{\unsure{g}}"']
\end{tikzcd}}
(avec des notations légèrement changées p.r. à (4.2))}
compte tenu que, par la structure \unsure{spéciale|précise} de $\pi=I$, on a
\[
\tag*{(6.2.5)}\mathbb{R}\Gamma^\pi(\Lambda_\nu)\simeq\Lambda_\nu\oplus\underbrace{\Lambda_\nu(-1)}_{\text{twist de Tate}}\overbrace{[-1]}^{\text{transl. des degrés}}.
\]
\deleted{L'homomorphisme $\mathbb{R}\Gamma^\pi(\kappa^*)$ \unsure{défini|définit} se décompose en
$\mathbb{R}\Gamma(V_\eta,\Lambda_\nu)\longrightarrow$}
Ainsi on retrouve $\mathbb{R}\Gamma(V_\eta,\Lambda_\nu)$, et l'hom $h^*$ déduit de $V_{\bar\eta}\xrightarrow{h}V_\eta$ n'est autre que l'hom naturel ``d'inclusion'' $\mathbb{R}\Gamma^\pi(K^\bullet)\longrightarrow\deleted{\unsure{\mathbb{R}\Gamma}}\{K^\bullet\}$. De plus, l'hom $\mathbb{R}\Gamma^\pi(\kappa^*)$ se décompose en
\[
\tag*{(6.2.6)}\mathbb{R}\Gamma^\pi(\kappa^*)=(\alpha,\beta):\mathbb{R}\Gamma(V_\eta,\Lambda_\nu)=\mathbb{R}\Gamma^\pi(K^\bullet)\longrightarrow\mathbb{R}\Gamma(V_s,\Lambda_\nu)+\mathbb{R}\Gamma(V_s,\Lambda_\nu)\unsure{[-1]|(-1)}
\]
